\documentclass[11pt]{article}
\usepackage{preamble}
\setlength{\parskip}{4pt}

\begin{document}

\daybanner{AP Statistics \textperiodcentered\ Unit 5 \textperiodcentered\ Topic 5.2 \textperiodcentered\ B: 2027-02-08 \textperiodcentered\ E: 2027-04-02 \textperiodcentered\ TEACHER KEY}{What Does This Correlation Support?}

\sectionbanner{CED TETHER}

{\small
\begin{itemize}[leftmargin=1.5em,itemsep=2pt,topsep=2pt]
  \item \textbf{Skill 2.C} --- Calculate summary statistics, relative positions of points within a distribution, correlation, and predicted response.
  \item \textbf{Skill 4.B} --- Interpret statistical calculations and findings to assign meaning or assess a claim.
  \item \textbf{EU DAT-1} --- Regression models may allow us to predict responses to changes in an explanatory variable.
  \item DAT-1.B Determine the correlation for a linear relationship. [Skill 2.C]
  \item DAT-1.B.1 The correlation, r, gives the direction and quantifies the strength of the linear association between two quantitative variables.  DAT-1.B.2 The correlation coefficient can be calculated by: r = (1/(n-1)) $\times$ $\sum$(($x_i$ - $\bar{x}$)/s\_x)((y\_i - $\bar{y}$)/s\_y). However, the most common way to determine r is by using technology.  DAT-1.B.3 A correlation coefficient close to 1 or -1 does not necessarily mean that a linear model is appropriate.
  \item DAT-1.C Interpret the correlation for a linear relationship. [Skill 4.B]
  \item DAT-1.C.1 The correlation, r, is unit-free, and always between -1 and 1, inclusive. A value of r = 0 indicates that there is no linear association. A value of r = 1 or r = -1 indicates that there is a perfect linear association.  DAT-1.C.2 A perceived or real relationship between two variables does not mean that changes in one variable cause changes in the other. That is, correlation does not necessarily imply causation.
\end{itemize}
}
\textbf{Essential question:} How should a correlation, its sensitivity, and the study design shape one recommendation?\par
\textbf{DOK-3 skill:} 4.B \textperiodcentered\ \textbf{FRQ pattern:} {\small\texttt{strong-\allowbreak{}correlation-\allowbreak{}driven-\allowbreak{}by-\allowbreak{}one-\allowbreak{}point-\allowbreak{}and-\allowbreak{}causation-\allowbreak{}overreach}}\par
\textbf{DOK rationale:} Part (c) weighs competing claims using the day's numerical or graphical evidence and explains a limitation or consequence; the judgment requires linked reasoning beyond a calculation.\par

\frameworkphaseheader{Do Now $\rightarrow$ video follow-along $\rightarrow$ finish and turn in}{1 $\rightarrow$ 3}{45}{%
\begin{itemize}[leftmargin=1.5em,itemsep=2pt,topsep=2pt]
  \item Show the board at the bell and collect a one-sentence first take without confirming a claim.
  \item Run the named video follow-along, then allow about 10 minutes for parts (a)--(c).
  \item Ask for evidence linked to a claim. Collect the sheet and hand-score part (c) with E/P/I.
\end{itemize}
}{%
\begin{itemize}[leftmargin=1.5em,itemsep=2pt,topsep=2pt]
  \item Commit to a proposal, then complete the follow-along.
  \item Finish (a)--(c), revise the first take, and turn in.
\end{itemize}
}{%
\begin{itemize}[leftmargin=1.5em,itemsep=2pt,topsep=2pt]
  \item Which number or visible feature supports your claim?
  \item Why does the other claim go beyond that evidence?
  \item What would you tell someone who chose the other claim?
\end{itemize}
}{Circulate and ask students to connect their choice to evidence. Score part (c) using the three required reasoning elements; do not require dropped extension work.}

\begin{callout}{\IconWarn\ WATCH FOR}{calloutred}
\begin{itemize}[leftmargin=1.5em,itemsep=2pt,topsep=2pt]
  \item A choice with copied numbers but no explanation of how they support it.
  \item Treating a model or average as a guaranteed individual outcome.
\end{itemize}
\end{callout}

\sectionbanner{SCORING PART (c) --- E / P / I}

\textbf{Expected elements}
\begin{itemize}[leftmargin=1.5em,itemsep=2pt,topsep=2pt]
  \item \textbf{reasoning-1} (required) --- Supports Tessa using both correlations and the separated point
  \item \textbf{reasoning-2} (required) --- Keeps the full result and clearly labels the comparison with one point removed
  \item \textbf{reasoning-3} (required) --- Uses the observational design to reject a causal visitor forecast
\end{itemize}

{\small\renewcommand{\arraystretch}{1.25}
\noindent\begin{tabular}{|p{0.5in}|p{5.9in}|}
\hline
\textbf{E} & Includes all three required reasoning elements above, with correct contextual evidence. \\ \hline
\textbf{P} & Includes at least one substantive correct reasoning link above, but omits or misstates another required element. \\ \hline
\textbf{I} & Includes no substantive correct reasoning link above; an unsupported choice or copied number alone is insufficient. \\ \hline
\end{tabular}}

\textbf{Common mistakes}
\begin{itemize}[leftmargin=1.5em,itemsep=2pt,topsep=2pt]
  \item Discarding the separated point without labeling the sensitivity analysis
  \item Ignoring the scatterplot after computing $r$
  \item Treating observational association as causation
\end{itemize}

\clearpage
\focusbanner{TODAY'S PROBLEM --- annotated}

Suppose a nonprofit records monthly social-media posts and average weekend visitors at eight community arts centers. The pairs are $(2,50)$, $(3,47)$, $(4,54)$, $(5,45)$, $(6,53)$, $(7,48)$, $(8,51)$, and $(20,90)$. Technology gives $r=0.922$ for all eight centers and $r=0.095$ for the first seven.\par\smallskip \textbf{Omar:} Use that strong correlation to forecast the attendance gained from additional posts; every observed center belongs in the analysis.\par \textbf{Tessa:} Report the full analysis, but also examine the seven-center cluster and how posting levels were chosen before recommending a campaign.\par
\begin{center}
\begin{tikzpicture}
\begin{axis}[
  width=0.53\linewidth,
  height=1.44in,
  xmin=0, xmax=21, ymin=40, ymax=95,
  xlabel={Social-media posts during the month}, ylabel={Average weekend visitors}, grid=major,
  tick label style={font=\small}, label style={font=\small}
]
\addplot[only marks, mark=*, mark size=1.4pt, color=dayblue] coordinates {
  (2,50)
  (3,47)
  (4,54)
  (5,45)
  (6,53)
  (7,48)
  (8,51)
  (20,90)
};
\end{axis}
\end{tikzpicture}
\end{center}
\begin{firsttakebox}{FIRST TAKE (5 min)}
Which proposal seems more defensible, Omar's or Tessa's? Cite one number, plot feature, or design fact.\par
\answer{Not graded; the finish must preserve valid evidence while bounding the claim.}
\end{firsttakebox}

\textit{Rules callout ``READ r WITH THE PLOT AND DESIGN (keep on the page)'' is printed on the student sheet.}\par\medskip

\dokbadge{1}~\textbf{(a)} What do the sign and size of the full-data $r=0.922$ tell you?\par
\answer{It summarizes a strong positive linear association across all eight centers: more posts tend to go with more visitors.}

\dokbadge{2}~\textbf{(b)} Compare the two given correlations. Describe the plot feature that could explain the difference.\par
\answer{The full correlation is $0.922$, but the first seven have $r=0.095$, close to zero. Seven points form a cloud with little linear pattern; $(20,90)$ stands apart at the upper right.}

\dokbadge{3}~\textbf{(c)} Which approach would you recommend? Use both correlations and the plot, explaining why the full result should still be reported. Explain why observing posting levels does not show that extra posts cause more visitors. Write 3--4 sentences.\par
\answer{I recommend Tessa: $r=0.922$ summarizes all eight centers, while $r=0.095$ and the seven-point cloud show how strongly the separated point affects that summary. Both results should be reported so a reader sees the observed point and its effect. Posting levels were observed, so other differences between centers could explain the relationship. Omar's proposed campaign could spend money on a visitor increase that these data do not establish.}

\begin{summaryexitbox}{EXIT}
One thing the video changed about my first take: \hrulefill
\end{summaryexitbox}

\end{document}
